% !TeX program = xelatex
\documentclass[11pt,a4paper]{article}

% -------------------------------------------------
% Packages
% -------------------------------------------------
\usepackage{kotex}
\usepackage{fontspec}
\usepackage{amsmath,amssymb,bm,mathtools}
\usepackage[left=24mm,right=24mm,top=23mm,bottom=25mm]{geometry}
\usepackage{booktabs,array}
\usepackage{enumitem}
\usepackage{setspace}
\usepackage{xcolor}
\usepackage{colortbl}
\usepackage{tikz}
\usepackage{pifont}
\usetikzlibrary{positioning,arrows.meta,calc,fit,decorations.pathreplacing}
\usepackage[most]{tcolorbox}
\usepackage{listings}
\usepackage{hyperref}
\usepackage{fancyhdr}

\setstretch{1.14}
\setlength{\parindent}{0pt}
\setlength{\parskip}{0.55em}

% -------------------------------------------------
% Colors and boxes
% -------------------------------------------------
\definecolor{deepblue}{RGB}{34,71,120}
\definecolor{softblue}{RGB}{239,246,253}
\definecolor{softgray}{RGB}{246,247,249}
\definecolor{softgreen}{RGB}{239,249,242}
\definecolor{softyellow}{RGB}{255,249,229}
\definecolor{codegray}{RGB}{248,248,248}

\newtcolorbox{keybox}[1][]{
  colback=softblue,
  colframe=deepblue,
  boxrule=0.7pt,
  arc=2mm,
  title=#1,
  fonttitle=\bfseries
}

\newtcolorbox{thinkbox}[1][]{
  colback=softyellow,
  colframe=orange!70!black,
  boxrule=0.6pt,
  arc=2mm,
  title=#1,
  fonttitle=\bfseries
}

\newtcolorbox{remarkbox}[1][]{
  colback=softgray,
  colframe=black!45,
  boxrule=0.5pt,
  arc=2mm,
  title=#1,
  fonttitle=\bfseries
}

\newtcolorbox{mlbox}[1][]{
  colback=softgreen,
  colframe=green!45!black,
  boxrule=0.6pt,
  arc=2mm,
  title=#1,
  fonttitle=\bfseries
}

% -------------------------------------------------
% Code style
% -------------------------------------------------
\lstdefinestyle{pythonstyle}{
  language=Python,
  backgroundcolor=\color{codegray},
  basicstyle=\ttfamily\small,
  keywordstyle=\bfseries,
  commentstyle=\itshape\color{black!60},
  stringstyle=\color{black!75},
  numbers=left,
  numberstyle=\tiny\color{black!55},
  numbersep=8pt,
  frame=single,
  rulecolor=\color{black!25},
  breaklines=true,
  showstringspaces=false,
  tabsize=4,
  columns=fullflexible,
  keepspaces=true
}
\lstset{style=pythonstyle}

% -------------------------------------------------
% Math commands
% -------------------------------------------------
\newcommand{\R}{\mathbb{R}}
\newcommand{\grad}{\nabla}
\newcommand{\vect}[1]{\bm{#1}}
\newcommand{\T}{^{\mathsf T}}
\newcommand{\loss}{\mathcal{L}}
\newcommand{\sigmoid}{\sigma}
\newcommand{\pd}[2]{\frac{\partial #1}{\partial #2}}
\newcommand{\softmax}{\operatorname{softmax}}

% -------------------------------------------------
% Header / footer
% -------------------------------------------------
\pagestyle{fancy}
\fancyhf{}
\lhead{Introduction to Deep Learning and Practice}
\rhead{Techniques for Training: Answers}
\cfoot{\thepage}
\renewcommand{\headrulewidth}{0.4pt}
\setlength{\headheight}{14pt}

\hypersetup{
  colorlinks=true,
  linkcolor=deepblue,
  urlcolor=deepblue
}

% =================================================

% source tag of a problem
\newcommand{\src}[1]{{\unskip\nobreak\hfil\penalty50\hskip1em\hbox{}\nobreak\hfill\mbox{\small\color{black!50}[#1]}}}
% a problem block
\newlist{problems}{enumerate}{1}
\setlist[problems]{label=\textbf{\arabic*.},leftmargin=2.2em,itemsep=0.2em}
% =================================================
% a solution block
\newtcolorbox{sol}{breakable,enhanced,colback=softgray,colframe=black!35,boxrule=0.4pt,arc=1mm,
  left=2mm,right=2mm,top=1mm,bottom=1mm,before upper={\textbf{풀이.}\ }}
\setlist[problems]{label=\textbf{\arabic*.},leftmargin=2.2em,itemsep=0.9em}
\begin{document}

\begin{center}
{\LARGE\bfseries 연습문제 정답과 풀이}\\[0.35em]
{\large 학습 관련 기술 $\cdot$ 딥러닝개론및연습 2026년 2학기}
\end{center}

강의노트 끝의 연습문제 17개의 정답과 풀이이다. 문제를 다시 적고 그 아래에 풀이를 쓴다.
코드를 실행하는 문제(5, 8, 11, 14, 17번)의 숫자는 모두 \texttt{연습문제\_정답/code/exercise\_runs.py}를
이 컴퓨터에서 실제로 실행한 결과이다.
초깃값과 미니배치 순서는 노트의 코드와 같은 씨앗(\texttt{np.random.seed(1)}, \texttt{np.random.seed(2)})으로 고정하였다.

\begin{problems}

% ---------------------------------------------------------------- 1
\item $f(x,y)=\frac{1}{20}x^2+y^2$에서 SGD가 발산하지 않는 학습률 $\alpha$의 범위를 구하라.
또 이 범위 안에서 $x$ 방향의 수렴이 가장 빠른 $\alpha$와 그때 $x$가 한 단계에 몇 배가 되는지 구하라.

\begin{sol}
$\nabla f=\bigl(\frac{x}{10},\ 2y\bigr)$이므로 SGD 갱신은 두 좌표가 따로따로이다.
\[
x_{k+1}=x_k-\alpha\frac{x_k}{10}=\Bigl(1-\frac{\alpha}{10}\Bigr)x_k,\qquad
y_{k+1}=y_k-2\alpha y_k=(1-2\alpha)\,y_k .
\]
각 좌표는 매 단계 같은 수를 곱하는 등비수열이므로, 곱하는 수의 절댓값이 $1$ 이하이면 발산하지 않는다.
\[
\Bigl|1-\frac{\alpha}{10}\Bigr|\le1\iff0\le\alpha\le20,\qquad
|1-2\alpha|\le1\iff0\le\alpha\le1 .
\]
둘 다 만족해야 하므로 $\boxed{0<\alpha\le1}$이다.
($\alpha=1$이면 $y_{k+1}=-y_k$로 $y$가 $\pm y_0$ 사이를 진동하며 수렴하지 않고, $0<\alpha<1$이면 두 좌표 모두 $0$으로 수렴한다.)
경사가 급한 $y$ 방향이 학습률의 상한을 정한다.

이 범위에서 $x$의 배율 $1-\frac{\alpha}{10}$은 $\alpha$가 클수록 작으므로 $\alpha=1$에서 가장 빠르고,
그때 $x$는 한 단계에 $\boxed{0.9}$배가 된다.
즉 $y$ 방향 때문에 학습률을 키울 수 없어서, 완만한 $x$ 방향은 아무리 빨라도 한 단계에 $10\%$씩밖에 줄지 않는다.
SGD가 이 함수에서 느린 이유이다.
\end{sol}

% ---------------------------------------------------------------- 2
\item 모멘텀에서 gradient가 매 단계 같은 값 $\bm g$로 일정하다고 하자.
$\bm v_0=\bm 0$일 때 $\bm v_k$를 구하고, $k\to\infty$일 때 $\bm v_k\to-\dfrac{\alpha}{1-\beta}\,\bm g$임을 보이라.

\begin{sol}
$\bm v_{k+1}=\beta\bm v_k-\alpha\bm g$를 차례로 쓰면
\[
\bm v_1=-\alpha\bm g,\quad
\bm v_2=-\alpha\bm g(1+\beta),\quad
\bm v_3=-\alpha\bm g(1+\beta+\beta^2),\quad\ldots
\]
이므로 (귀납법으로 확인된다)
\[
\bm v_k=-\alpha\bm g\,(1+\beta+\cdots+\beta^{k-1})=-\alpha\,\frac{1-\beta^k}{1-\beta}\,\bm g .
\]
$0\le\beta<1$이면 $\beta^k\to0$이므로 $\bm v_k\to-\dfrac{\alpha}{1-\beta}\bm g$이다.

$\beta=0.9$이면 $\frac{1}{1-\beta}=10$이므로, gradient의 방향이 일정한 곳에서 모멘텀은 SGD보다 최대 $10$배 빠르게 움직인다.
반대로 gradient의 부호가 번갈아 바뀌는 방향에서는 $\bm v$에 쌓인 값이 서로 상쇄되어 진동이 줄어든다.
\end{sol}

% ---------------------------------------------------------------- 3
\item 시작점 $(-7,2)$에서 AdaGrad ($\alpha=1.5$)로 두 번째 단계 $(x_2,y_2)$를 손으로 계산하라
($\varepsilon$은 무시한다). 표의 값 $(-4.573,\ 0.136)$과 비교하라.

\begin{sol}
갱신식은 $\bm h\leftarrow\bm h+\bm g\odot\bm g$, $W\leftarrow W-\alpha\,\bm g/\sqrt{\bm h}$ (성분별)이고 $\nabla f=(x/10,\ 2y)$이다.

\textbf{1단계.} $\bm g_0=(-0.7,\ 4)$, $\bm h_1=(0.49,\ 16)$.
\[
x_1=-7-1.5\cdot\frac{-0.7}{\sqrt{0.49}}=-7+1.5=-5.5,\qquad
y_1=2-1.5\cdot\frac{4}{\sqrt{16}}=2-1.5=0.5 .
\]
첫 단계에서는 $\sqrt{\bm h_1}=|\bm g_0|$이므로 두 좌표가 모두 정확히 $\alpha=1.5$만큼 움직인다.

\textbf{2단계.} $\bm g_1=(-0.55,\ 1)$, $\bm h_2=(0.49+0.3025,\ 16+1)=(0.7925,\ 17)$.
\[
x_2=-5.5+1.5\cdot\frac{0.55}{\sqrt{0.7925}}=-5.5+\frac{0.825}{0.8902}\approx-5.5+0.9267=-4.573,
\]
\[
y_2=0.5-1.5\cdot\frac{1}{\sqrt{17}}=0.5-\frac{1.5}{4.1231}\approx0.5-0.3638=0.136 .
\]
$\boxed{(x_2,y_2)\approx(-4.573,\ 0.136)}$으로 표의 값과 같다.
$y$ 방향은 첫 단계의 큰 gradient ($4$)가 $h$에 쌓여 있어서, $\bm g$가 $1$인데도 $0.36$만 움직인다.
$x$ 방향은 $h$가 작아서 gradient $0.55$에 비해 크게 ($0.93$) 움직인다.
\end{sol}

% ---------------------------------------------------------------- 4
\item Adam에서 편향 보정을 하지 않으면 첫 단계의 이동량은 얼마가 되는가?
($\beta_1=0.9$, $\beta_2=0.999$, $\varepsilon$은 무시한다.) 편향 보정이 왜 필요한지 설명하라.

\begin{sol}
$\bm m_0=\bm v_0=\bm 0$에서 첫 gradient $g$ (성분 하나)를 받으면
\[
m_1=(1-\beta_1)g=0.1\,g,\qquad v_1=(1-\beta_2)g^2=0.001\,g^2 .
\]
보정하지 않으면 이동량은
\[
\alpha\frac{m_1}{\sqrt{v_1}}=\alpha\frac{0.1\,g}{\sqrt{0.001}\,|g|}
=\frac{0.1}{0.0316}\,\alpha\operatorname{sign}(g)\approx\boxed{3.16\,\alpha}
\]
이다 (방향은 $-\operatorname{sign}(g)$).
보정하면 $\hat m_1=m_1/(1-\beta_1)=g$, $\hat v_1=v_1/(1-\beta_2)=g^2$이므로 이동량은 정확히 $\alpha$이다.

\textbf{편향 보정이 필요한 이유.} $\bm m$, $\bm v$는 $\bm 0$에서 시작하는 지수이동평균이므로 처음 몇 단계에서는 $0$ 쪽으로 치우쳐 있다
(gradient가 일정하면 $m_k=(1-\beta_1^k)g$, $v_k=(1-\beta_2^k)g^2$).
$\beta_2=0.999$가 $\beta_1=0.9$보다 $1$에 훨씬 가까워서 $v$가 훨씬 더 오래 작게 머문다.
그래서 분모 $\sqrt{v}$가 너무 작아지고, 학습 초기의 이동량이 의도한 $\alpha$보다 커진다 (첫 단계에서 약 $3$배).
$1-\beta^k$로 나누면 $\hat m_k$, $\hat v_k$가 처음부터 $g$, $g^2$의 크기를 갖게 되어, 학습률 $\alpha$가 처음부터 실제 이동량의 크기를 뜻하게 된다.
$k$가 커지면 $1-\beta^k\to1$이므로 보정의 효과는 사라진다.
\end{sol}

% ---------------------------------------------------------------- 5
\item \texttt{mnist\_optimizer\_compare.py}에서 SGD의 학습률을 $0.01$에서 $0.05$, $0.1$로 바꾸어 보라.
SGD도 plateau를 빠져나오는가? 같은 조건에서 RMSProp ($\alpha=0.001$)을 \texttt{optimizers}에 추가하여
네 기법과 비교하라.

\begin{sol}
노트와 같은 조건 (은닉층 $[100]\times4$, 초깃값 표준편차 $0.05$, 미니배치 $128$, $5$ epoch, 같은 초깃값과 같은 미니배치)에서 실행한 검증 정확도이다.
\begin{center}
\begin{tabular}{lccccc}
\toprule
기법 & epoch 1 & 2 & 3 & 4 & 5 \\
\midrule
SGD ($\alpha=0.01$) & 0.1138 & 0.1144 & 0.2046 & 0.2824 & 0.4404 \\
SGD ($\alpha=0.05$) & 0.4210 & 0.8404 & 0.9098 & 0.9274 & 0.9166 \\
SGD ($\alpha=0.1$)  & 0.7802 & 0.9134 & 0.9512 & 0.9584 & 0.9584 \\
\midrule
Momentum ($\alpha=0.01$, $\beta=0.9$) & 0.8448 & 0.9286 & 0.9464 & 0.9588 & 0.9632 \\
AdaGrad ($\alpha=0.01$) & 0.9448 & 0.9602 & 0.9660 & 0.9688 & \textbf{0.9702} \\
RMSProp ($\alpha=0.001$, $\rho=0.9$) & 0.9312 & 0.9498 & 0.9662 & 0.9690 & 0.9700 \\
Adam ($\alpha=0.001$) & 0.9264 & 0.9548 & 0.9586 & 0.9668 & 0.9660 \\
\bottomrule
\end{tabular}
\end{center}
\begin{itemize}[leftmargin=1.5em,itemsep=2pt]
\item \textbf{SGD도 plateau를 빠져나온다.} $\alpha=0.01$에서는 $2$ epoch 동안 정확도가 $11\%$ (찍는 수준)에 머물렀지만,
$\alpha=0.05$에서는 $2$ epoch 만에 $84\%$, $\alpha=0.1$에서는 $1$ epoch 만에 $78\%$가 되었다.
plateau는 gradient가 아주 작은 평평한 영역이므로, 학습률을 키우면 같은 작은 gradient로도 더 멀리 움직여 빨리 빠져나온다.
$\alpha=0.1$의 SGD는 $5$ epoch 후 $95.84\%$로 Momentum ($96.32\%$)에 가깝다.
즉 노트의 비교에서 SGD가 크게 뒤진 것은 상당 부분 학습률 $0.01$이 SGD에게 너무 작았기 때문이다.
\item $\alpha=0.05$에서 epoch 4에서 5로 갈 때 정확도가 조금 떨어졌다 ($0.9274\to0.9166$). 미니배치 gradient의 잡음 때문이며, epoch마다 값이 오르내릴 수 있다.
\item \textbf{RMSProp}은 $97.00\%$로 AdaGrad ($97.02\%$)와 거의 같고 Adam ($96.60\%$)보다 조금 높다.
첫 epoch부터 $93\%$로, 적응적 학습률을 쓰는 세 기법은 모두 plateau 없이 바로 학습을 시작한다.
다만 $5$ epoch, 씨앗 하나의 결과이므로 $0.1\%$p 수준의 차이로 순위를 매기기는 어렵다.
\end{itemize}
RMSProp은 \texttt{optimizers.py}에 이미 있는 \texttt{RMSProp(lr=0.001)}을 \texttt{optimizers} 사전에 한 줄 추가하면 된다.
\end{sol}

% ---------------------------------------------------------------- 6
\item 은닉층의 모든 가중치를 같은 값 $c$로 초기화하면 학습이 진행되어도 은닉층의 node들이
계속 같은 값을 가지는 이유를 역전파 식 $\pd{L}{W}=X\T\pd{L}{A}$를 이용하여 설명하라.

\begin{sol}
2층 network $A=XW^{(1)}+B^{(1)}$, $Z=h(A)$, $U=ZW^{(2)}+B^{(2)}$로 설명한다.
$W^{(1)}$의 모든 성분이 $c$이고 $B^{(1)}$의 성분도 모두 같으면, $W^{(1)}$의 열이 모두 같다.

\textbf{순전파.} $A$의 $j$번째 열은 $XW^{(1)}_{:,j}+b_j$이므로 모든 열이 같다. 따라서 $Z$의 열 (은닉 node의 값)도 모두 같다.

\textbf{역전파.} $W^{(2)}$의 모든 성분도 $c$이면 $W^{(2)}$의 행이 모두 같으므로 $\pd{L}{Z}=\pd{L}{U}\,W^{(2)\mathsf T}$의 열이 모두 같다.
$Z$의 열이 같으므로 $h'(A)$의 열도 같고, $\pd{L}{A}=\pd{L}{Z}\odot h'(A)$의 열도 모두 같다.
그러면
\[
\pd{L}{W^{(1)}}=X\T\pd{L}{A}
\]
의 $j$번째 열은 $X\T\bigl(\pd{L}{A}\bigr)_{:,j}$이므로 \textbf{모든 열이 같다.} 마찬가지로 $\pd{L}{B^{(1)}}$의 성분도 모두 같다.
또 $\pd{L}{W^{(2)}}=Z\T\pd{L}{U}$의 행은 $Z$의 열로 정해지므로 모든 행이 같다.

따라서 갱신 $W\leftarrow W-\alpha\pd{L}{W}$ 후에도 $W^{(1)}$의 열은 서로 같고 $W^{(2)}$의 행도 서로 같다.
같은 논리가 매 단계 반복되므로 은닉 node들은 학습 내내 같은 값을 가진다.
은닉층에 node가 $100$개 있어도 node $1$개와 같은 표현력밖에 없다.
초깃값을 무작위로 정하는 것은 이 대칭을 깨기 위해서이다.
($c=0$이면 $\pd{L}{Z}=0$이 되어 $W^{(1)}$은 아예 갱신되지 않는다.)
\end{sol}

% ---------------------------------------------------------------- 7
\item $a=\sum_{i=1}^{n}z_iw_i$에서 $z_i=\max(0,a_i')$이고 $a_i'$이 평균 $0$인 대칭 분포를 따를 때
$\mathbb E[z_i^2]=\frac12\mathrm{Var}(a_i')$임을 보이고, 이로부터 He 초깃값 $\mathrm{Var}(w)=\frac2n$을 유도하라.

\begin{sol}
$a_i'$의 확률밀도를 $p$라 하면 $p(-s)=p(s)$이다. $s<0$이면 $z_i=0$이므로
\[
\mathbb E[z_i^2]=\int_0^{\infty}s^2p(s)\,ds
=\frac12\int_{-\infty}^{\infty}s^2p(s)\,ds
=\frac12\mathbb E[(a_i')^2]=\frac12\mathrm{Var}(a_i') .
\]
(둘째 등호는 $s^2p(s)$가 짝함수이기 때문이고, 마지막 등호는 평균이 $0$이기 때문이다.)

$w_i$는 $z_i$와 독립이고 평균 $0$, 분산 $\mathrm{Var}(w)$이며, 항들이 서로 독립이라 하자.
$\mathbb E[z_iw_i]=\mathbb E[z_i]\mathbb E[w_i]=0$이므로
\[
\mathrm{Var}(a)=\sum_{i=1}^{n}\mathrm{Var}(z_iw_i)=\sum_{i=1}^{n}\mathbb E[z_i^2]\,\mathbb E[w_i^2]
=n\cdot\frac12\mathrm{Var}(a')\cdot\mathrm{Var}(w) .
\]
($z_i$의 평균은 $0$이 아니므로 $\mathrm{Var}(z_i)$가 아니라 $\mathbb E[z_i^2]$이 들어간다는 점에 주의한다.)
층을 지나도 활성화 전 값의 분산이 유지되려면 ($\mathrm{Var}(a)=\mathrm{Var}(a')$)
\[
\frac n2\,\mathrm{Var}(w)=1\quad\Longrightarrow\quad\boxed{\mathrm{Var}(w)=\frac2n},\qquad\text{즉 표준편차 }\sqrt{2/n}.
\]
Xavier 초깃값 $\frac1n$과 비교하면, ReLU가 음수 절반을 $0$으로 만들어 신호의 크기가 절반으로 줄어드는 것을 $2$배로 보상한 것이다.
\end{sol}

% ---------------------------------------------------------------- 8
\item \texttt{weight\_init\_activation.py}에서 층의 수를 $5$에서 $20$으로 늘려
ReLU + Xavier와 ReLU + He의 마지막 층 활성화값의 표준편차를 비교하라.
Xavier에서 표준편차가 층마다 약 몇 배씩 줄어드는가?

\begin{sol}
노트의 코드와 같은 조건 (입력 $1000\times100$, node $100$개, bias $0$, \texttt{default\_rng(0)})에서 층을 $20$개로 늘린 결과이다.
처음 $5$층의 값은 노트의 결과와 같다.
\begin{center}
\begin{tabular}{lcccccc}
\toprule
층 & 1 & 5 & 10 & 15 & 20 & 층당 배율 (기하평균) \\
\midrule
ReLU + Xavier & 0.583 & 0.177 & 0.0327 & 0.00812 & \textbf{0.00198} & 0.741 \\
ReLU + He     & 0.825 & 1.000 & 1.047 & 1.469 & \textbf{2.025} & 1.048 \\
\bottomrule
\end{tabular}
\end{center}
마지막 ($20$번째) 층의 표준편차는 Xavier가 $0.00198$, He가 $2.03$으로 약 $1000$배 차이가 난다.
Xavier에서는 표준편차가 층마다 약 $\boxed{0.74\text{배}}$씩 줄어든다.

이론값은 $\frac{1}{\sqrt2}\approx0.71$배이다.
문제 7의 식에서 $\mathrm{Var}(w)=\frac1n$이면 $\mathrm{Var}(a)=\frac n2\cdot\frac1n\mathrm{Var}(a')=\frac12\mathrm{Var}(a')$이므로
분산이 층마다 절반, 표준편차가 $\frac1{\sqrt2}$배가 된다. $20$층을 지나면 $(1/\sqrt2)^{19}\approx0.0014$배가 된다.
실험값이 조금 다른 것은 node가 $100$개뿐이라 무작위 가중치에 따른 변동이 있기 때문이다.
He는 이론적으로 층마다 $1$배이며, 실험에서도 크게 줄거나 늘지 않는다
(층당 $1.05$배로, $20$층 동안 $0.8$에서 $2.0$ 사이에 머문다).
층이 깊을수록 초깃값의 작은 차이가 거듭제곱으로 커지므로 ReLU에는 He 초깃값을 써야 한다.
\end{sol}

% ---------------------------------------------------------------- 9
\item 배치 정규화의 순전파 식 (미니배치 평균, 분산, 정규화, 확대와 이동)을 보고,
같은 미니배치에서 Affine 계층의 가중치 $W$를 $a$배 ($a>0$) 하면 BatchNorm 계층의 출력 $Y$가 변하지 않음을 보이라
($\varepsilon$은 무시한다). 이 사실로부터 $\pd{L}{(aW)}=\frac1a\pd{L}{W}$임을 설명하라.

\begin{sol}
BatchNorm의 입력을 $x_i=x_{\text{in},i}W+B$ (node 하나의 성분)라 하고, $W$를 $aW$로 바꾼 값을 $x_i'$라 하자.
$x_i'=a\,x_{\text{in},i}W+B$이므로
\[
\mu'=a\,\overline{x_{\text{in}}W}+B,\qquad
x_i'-\mu'=a\bigl(x_{\text{in},i}W-\overline{x_{\text{in}}W}\bigr)=a\,(x_i-\mu).
\]
($B$는 모든 행에 같은 값이므로 평균을 빼면 사라진다.) 따라서
\[
\sigma'^2=\frac1h\sum_i(x_i'-\mu')^2=a^2\sigma^2,\qquad
\hat x_i'=\frac{x_i'-\mu'}{\sqrt{\sigma'^2}}=\frac{a(x_i-\mu)}{a\sigma}=\hat x_i,
\]
이고 $y_i'=\gamma\hat x_i'+\beta=y_i$이다. 즉 $W$를 몇 배 하든 BatchNorm의 출력, 그리고 그 뒤의 손실 $L$도 그대로이다.

\textbf{gradient.} 위의 사실은 모든 $W$에 대하여 $L(aW)=L(W)$라는 뜻이다.
양변을 $W$로 미분하면 왼쪽은 연쇄법칙으로 $a\cdot\nabla L(aW)$이므로
\[
a\,\nabla L(aW)=\nabla L(W)\qquad\Longrightarrow\qquad
\pd{L}{(aW)}=\frac1a\pd{L}{W}.
\]
(여기서 $\pd{L}{(aW)}$는 가중치가 $aW$인 점에서 구한 gradient이다.)
가중치가 커지면 ($a>1$) gradient가 그만큼 작아지고, 작아지면 gradient가 커진다.
그래서 가중치의 크기가 저절로 조절되어, BatchNorm을 쓰면 초깃값의 크기나 학습률에 덜 민감해진다.
노트의 실험에서 초깃값의 표준편차를 바꾸어도 BatchNorm을 쓴 network가 잘 학습한 것과 같은 이야기이다.
\end{sol}

% ---------------------------------------------------------------- 10
\item 배치 정규화 계산 그래프의 역전파 단계 \ding{173}--\ding{178}을 모두 대입하여
$\pd{L}{x_i}=\frac{1}{\sqrt{\sigma_B^2+\varepsilon}}\bigl(g_i-\bar g-\hat x_i\,\overline{g\hat x}\bigr)$을 유도하라.
($g_i=\pd{L}{\hat x_i}$이고, $\bar g$, $\overline{g\hat x}$는 각각 $g_k$, $g_k\hat x_k$의 미니배치 평균이다.)

\begin{sol}
node 하나에 대하여 쓴다. $s=\sqrt{\sigma_B^2+\varepsilon}$, $x_{c,i}=x_i-\mu$, $\hat x_i=x_{c,i}/s$이다.
\begin{itemize}[leftmargin=2.2em,itemsep=4pt]
\item[\ding{173}] $g_i=\pd{L}{\hat x_i}=\gamma\pd{L}{y_i}$.
\item[\ding{174}] $\displaystyle\pd{L}{s}=-\sum_k g_k\frac{x_{c,k}}{s^2}$.
\item[\ding{175}] $\displaystyle\pd{L}{v}=\frac{1}{2s}\pd{L}{s}=-\frac{1}{2s^3}\sum_k g_kx_{c,k}$.
\item[\ding{176}] $\displaystyle\pd{L}{x_{c,i}}=\frac{g_i}{s}+\frac2h\,x_{c,i}\pd{L}{v}
=\frac{g_i}{s}-\frac{x_{c,i}}{h\,s^3}\sum_kg_kx_{c,k}.$\\[2pt]
$x_{c,i}=s\hat x_i$를 넣으면 $\displaystyle\frac{x_{c,i}}{h s^3}\sum_k g_kx_{c,k}=\frac{\hat x_i}{s}\cdot\frac1h\sum_kg_k\hat x_k=\frac{\hat x_i}{s}\,\overline{g\hat x}$이므로
\[
\pd{L}{x_{c,i}}=\frac1s\bigl(g_i-\hat x_i\,\overline{g\hat x}\bigr).
\]
\item[\ding{177}] $\displaystyle\pd{L}{\mu}=-\sum_k\pd{L}{x_{c,k}}=-\frac1s\Bigl(\sum_kg_k-\overline{g\hat x}\sum_k\hat x_k\Bigr)=-\frac hs\,\bar g$.\\[2pt]
여기서 $\hat x$의 미니배치 평균이 $0$, 즉 $\sum_k\hat x_k=0$임을 썼다.
\item[\ding{178}] $\displaystyle\pd{L}{x_i}=\pd{L}{x_{c,i}}+\frac1h\pd{L}{\mu}
=\frac1s\bigl(g_i-\hat x_i\,\overline{g\hat x}\bigr)-\frac1s\bar g .$
\end{itemize}
따라서
\[
\boxed{\pd{L}{x_i}=\frac{1}{\sqrt{\sigma_B^2+\varepsilon}}\bigl(g_i-\bar g-\hat x_i\,\overline{g\hat x}\bigr)}.
\]
$\sum_i\hat x_i=0$이고 $\sum_i\hat x_i^2=h$ ($\varepsilon$ 무시)이므로
$\sum_i\pd{L}{x_i}=\frac1s(h\bar g-h\bar g-0)=0$이고 $\sum_i\hat x_i\pd{L}{x_i}=\frac1s(h\overline{g\hat x}-0-h\overline{g\hat x})=0$이다.
즉 입력 gradient에는 ``모든 데이터에 같은 값을 더하는 방향''과 ``$\hat x$ 방향으로 늘리는 방향'' 성분이 없다.
정규화가 그 두 변화를 지워 버리기 때문이다.
\end{sol}

% ---------------------------------------------------------------- 11
\item \texttt{mnist\_batchnorm.py}의 실험 2에서 \texttt{accuracy}가 이동평균 대신 미니배치 통계를 쓰도록
(\texttt{train\_flg=True}) 바꾸고, 검증 데이터를 $1$장씩 넣어 정확도를 계산해 보라.
어떤 일이 생기는가? 왜 그런가?

\begin{sol}
실험 2의 조건 (훈련 이미지 $1000$장, 은닉층 $[100]\times5$, BatchNorm, SGD $\alpha=0.01$, $20$ epoch) 중 초깃값 표준편차 $0.1$인 network를 학습한 뒤 검증 데이터 $5000$장으로 정확도를 계산하였다.
\begin{center}
\begin{tabular}{lc}
\toprule
추론에 쓴 통계 & 검증 정확도 \\
\midrule
이동평균 (\texttt{train\_flg=False}, 원래 코드) & 0.7030 \\
미니배치 통계, $5000$장을 한꺼번에 & 0.7038 \\
미니배치 통계, $1$장씩 & \textbf{0.1000} \\
\bottomrule
\end{tabular}
\end{center}
$1$장씩 넣으면 \textbf{$5000$장 모두를 같은 class (9)로 예측}하여 정확도가 $10\%$ (검증 데이터에서 9의 비율)로 떨어진다.
두 이미지의 점수 $U$를 비교하면 완전히 같다.

\textbf{이유.} 미니배치 크기가 $h=1$이면 평균은 그 데이터 자신 ($\mu=x$)이므로
\[
x-\mu=0,\qquad \sigma^2=0,\qquad \hat x=\frac{0}{\sqrt{0+\varepsilon}}=0,\qquad y=\gamma\cdot0+\beta=\beta .
\]
첫 BatchNorm 계층의 출력이 입력과 관계없이 항상 $\beta$가 되므로, 그 뒤의 모든 계층의 값도 입력과 관계없이 같다.
즉 입력 이미지의 정보가 첫 BatchNorm에서 모두 지워진다.
$5000$장을 한꺼번에 넣으면 그 통계가 이동평균과 비슷하므로 정확도가 거의 같다.
하지만 그러면 한 이미지의 예측이 함께 넣은 다른 이미지에 따라 달라진다.

그래서 추론할 때는 학습 중에 모아 둔 이동평균 (\texttt{running\_mean}, \texttt{running\_var})을 써야 한다.
그러면 데이터 한 장도 다른 데이터와 상관없이 예측할 수 있다.
PyTorch에서는 \texttt{model.eval()}이 이 전환을 한다.
\end{sol}

% ---------------------------------------------------------------- 12
\item 가중치 감소를 적용한 SGD의 갱신식이 $W\leftarrow(1-\alpha\lambda)W-\alpha\pd{L}{W}$임을 보이라.
$\alpha=0.01$, $\lambda=0.1$일 때 gradient가 $0$이라면 $W$가 처음 값의 절반이 되기까지 몇 번의 갱신이 필요한가?

\begin{sol}
가중치 감소를 넣은 손실은 $\tilde L=L+\frac\lambda2\|W\|^2$이고 $\pd{\tilde L}{W}=\pd{L}{W}+\lambda W$이다. SGD로 갱신하면
\[
W\leftarrow W-\alpha\Bigl(\pd{L}{W}+\lambda W\Bigr)=(1-\alpha\lambda)W-\alpha\pd{L}{W}.
\]
gradient가 $0$이면 매번 $W$에 $1-\alpha\lambda=1-0.001=0.999$를 곱한다.
$k$번 갱신 후 $W_k=0.999^kW_0$이므로 $0.999^k\le\frac12$인 가장 작은 $k$를 구한다.
\[
k\ge\frac{\ln\frac12}{\ln0.999}=\frac{-0.693147}{-0.0010005}\approx692.8
\quad\Longrightarrow\quad\boxed{693\text{번}}.
\]
($\ln(1-x)\approx-x$를 쓰면 $k\approx\frac{\ln2}{\alpha\lambda}=693$으로 같은 답을 얻는다.)
MNIST에서 미니배치 $100$이면 $1$ epoch가 $550$번의 갱신이므로, 약 $1.3$ epoch 동안 gradient가 없는 가중치는 절반이 된다.
\end{sol}

% ---------------------------------------------------------------- 13
\item 하나의 parameter $w$에 대하여 gradient가 매 단계 $g$로 일정하다고 하자 ($g\neq0$, $\varepsilon$은 무시).
Adam에 L2 규제를 쓰면 감소 항 $\lambda w$가 실제 갱신량에 주는 영향이 $|g|$에 따라 어떻게 달라지는지,
AdamW와 비교하여 설명하라.

\begin{sol}
\textbf{Adam + L2.} L2 규제는 손실에 $\frac\lambda2w^2$을 더하므로 Adam이 받는 gradient는 $g'=g+\lambda w$이다.
gradient가 일정하면 (편향 보정 후) $\hat m=g'$, $\hat v=g'^2$이므로 갱신량은
\[
\Delta w=-\alpha\frac{\hat m}{\sqrt{\hat v}}=-\alpha\frac{g+\lambda w}{|g+\lambda w|}=-\alpha\operatorname{sign}(g+\lambda w)
\]
이다. 갱신량의 크기는 항상 $\alpha$이고, 감소 항 $\lambda w$는 \textbf{분모 $\sqrt{\hat v}$로 함께 나누어진다}.
\begin{itemize}[leftmargin=1.5em,itemsep=2pt]
\item $|g|\gg\lambda|w|$이면 $\operatorname{sign}(g+\lambda w)=\operatorname{sign}(g)$이므로 감소 항이 갱신량을 거의 바꾸지 못한다.
감소 항의 효과는 대략 $\alpha\lambda w/|g|$ 정도로, $|g|$에 반비례하여 작아진다.
\item $|g|\ll\lambda|w|$이면 갱신이 $-\alpha\operatorname{sign}(w)$가 되어, $w$가 매 단계 $\alpha$씩 $0$으로 끌려간다.
\end{itemize}
즉 gradient가 큰 parameter일수록 규제를 덜 받는다. 규제의 세기가 parameter마다 다르게, 의도와 상관없이 정해진다.

\textbf{AdamW.} 감소를 Adam의 계산과 분리하여 따로 한다.
\[
w\leftarrow w-\alpha\lambda w-\alpha\frac{\hat m}{\sqrt{\hat v}}=w-\alpha\lambda w-\alpha\operatorname{sign}(g).
\]
감소량 $\alpha\lambda w$는 $g$와 관계없이 모든 parameter에서 같은 비율이다.
SGD에서 L2 규제와 가중치 감소가 같았던 성질 (문제 12)을 Adam에서도 되살린 것이다.
노트의 실험 (훈련 이미지 $300$장)에서 Adam + L2 ($\lambda=0.1$)는 시험 정확도가 $72.90\%$로 떨어졌지만 AdamW ($0.1$)는 $83.12\%$였던 것도 이 차이에서 나온다.
\end{sol}

% ---------------------------------------------------------------- 14
\item \texttt{mnist\_overfit.py}에서 은닉층을 $[100]\times6$ 대신 $[500, 500]$으로 바꾸고,
드롭아웃의 삭제 비율 $p=0,\ 0.2,\ 0.5$를 \textbf{검증 데이터}로 비교하라.
어떤 $p$를 고르겠는가? 그 $p$로 학습한 network의 시험 정확도를 보고하라.

\begin{sol}
노트의 조건 (훈련 이미지 $300$장, He 초깃값, SGD $\alpha=0.01$, 미니배치 $100$, $300$ epoch, 같은 초깃값과 미니배치 순서)에서 은닉층만 $[500,500]$으로 바꾸었다.
\begin{center}
\begin{tabular}{cccc}
\toprule
$p$ & 훈련 정확도 & 검증 정확도 & 시험 정확도 \\
\midrule
$0$   & 1.0000 & \textbf{0.8314} & (0.8309) \\
$0.2$ & 1.0000 & 0.8286 & ---\\
$0.5$ & 0.9800 & 0.8206 & ---\\
\bottomrule
\end{tabular}
\end{center}
검증 정확도가 가장 높은 $\boxed{p=0}$ (드롭아웃 없음)을 고른다.
그 network의 시험 정확도는 $\boxed{83.09\%}$이다.
시험 데이터는 $p$를 고른 \textbf{뒤에} 한 번만 계산하였고, 나머지 두 $p$의 시험 정확도는 계산하지 않았다
(계산해서 비교하면 시험 데이터로 고르는 것이 된다).

\textbf{해석.}
\begin{itemize}[leftmargin=1.5em,itemsep=2pt]
\item 이 설정에서는 드롭아웃이 도움이 되지 않았다. 세 값의 차이는 $1\%$p 이내 ($p=0$과 $0.2$의 차이 $0.28\%$p는 검증 이미지 $14$장)로 작다.
층이 $2$개뿐인 넓은 network는 깊은 $[100]\times6$보다 덜 불안정하고, 학습률 $0.01$의 SGD로는 드롭아웃의 잡음이 학습을 늦추는 효과가 더 컸다고 볼 수 있다.
\item 훈련 정확도는 모두 $100\%$ 근처이므로 오버피팅은 여전하다. 훈련 데이터가 $300$장뿐이라는 근본 문제는 규제만으로 해결되지 않는다.
\item 검증 이미지 $5000$장에서 생기는 정확도의 표준오차는 약 $\sqrt{0.83\cdot0.17/5000}\approx0.5\%$p이므로, 이 정도 차이로 ``드롭아웃은 해롭다''고 결론 내릴 수는 없다. 씨앗을 바꾸어 여러 번 반복하면 더 믿을 만한 비교가 된다.
\end{itemize}
\end{sol}

% ---------------------------------------------------------------- 15
\item 하이퍼파라미터 $10$개를 시험하여 시험 정확도가 가장 높은 것을 골랐더니 시험 정확도가 $98\%$였다.
이 $98\%$를 처음 보는 데이터에 대한 성능이라고 보고해도 되는가? 왜 그런가? 어떻게 해야 하는가?

\begin{sol}
\textbf{보고하면 안 된다.}
시험 데이터로 $10$개 중 하나를 \textbf{골랐으므로} 시험 데이터가 학습 과정 (모델 선택)에 쓰였다.
$10$개의 시험 정확도에는 각각 우연한 오차가 있는데, 그중 최댓값을 고르면 우연히 운이 좋았던 것을 고르게 된다.
그래서 $98\%$는 처음 보는 데이터에 대한 성능보다 높게 (낙관적으로) 추정된 값이다.
하이퍼파라미터가 시험 데이터에 오버피팅된 것이라고 할 수 있다.

\textbf{올바른 방법.}
\begin{enumerate}[label=(\arabic*),leftmargin=2em,itemsep=1pt]
\item 훈련 데이터의 일부를 \textbf{검증 데이터}로 떼어 놓는다 (노트에서는 $60000$장 중 $5000$장).
\item $10$개의 하이퍼파라미터를 훈련 데이터로 학습하고 \textbf{검증 정확도}로 비교하여 하나를 고른다.
\item 고른 network의 시험 정확도를 \textbf{마지막에 한 번만} 계산하여 보고한다.
\end{enumerate}
이미 시험 데이터로 골랐다면, 그 결과를 보고하는 대신 고른 하이퍼파라미터를 아직 쓰지 않은 새 데이터로 평가해야 한다.
\end{sol}

% ---------------------------------------------------------------- 16
\item 학습률 $\alpha$를 $[10^{-5},\,1]$에서 (a) 그냥 균등하게, (b) $\alpha=10^{u}$, $u\sim\mathrm{Uniform}(-5,0)$으로 뽑을 때,
$\alpha<10^{-2}$일 확률을 각각 구하라.

\begin{sol}
(a) $\alpha\sim\mathrm{Uniform}(10^{-5},1)$이면
\[
P(\alpha<10^{-2})=\frac{10^{-2}-10^{-5}}{1-10^{-5}}=\frac{0.00999}{0.99999}\approx\boxed{0.0100}.
\]
(b) $\alpha<10^{-2}\iff u<-2$이므로
\[
P(u<-2)=\frac{-2-(-5)}{0-(-5)}=\boxed{\frac35=0.6}.
\]
(a)로 뽑으면 $99\%$가 $[0.01,\,1]$에 몰리고, $[10^{-5},10^{-4}]$, $[10^{-4},10^{-3}]$처럼 작은 구간은 거의 뽑히지 않는다.
(b)는 $[10^{-5},10^{-4}]$, $\ldots$, $[10^{-1},1]$의 다섯 구간 (자릿수)을 각각 $20\%$씩 고르게 뽑는다.
학습률의 효과는 값의 차이가 아니라 \textbf{비율} (몇 배인가)에 따라 달라지므로 로그 척도로 뽑는다.
\end{sol}

% ---------------------------------------------------------------- 17
\item \texttt{hyperparameter\_search.py}의 결과를 보고 학습률과 가중치 감소의 범위를 좁혀
($100$번 대신 $30$번) 다시 탐색하라. 좁힌 범위와 그 근거, 가장 좋은 검증 정확도를 표로 정리하라.

\begin{sol}
\textbf{원래 탐색 ($100$번)의 결과} (노트의 결과를 다시 실행한 값).
\begin{itemize}[leftmargin=1.5em,itemsep=1pt]
\item 학습률: 검증 정확도 상위 $10$개가 모두 $\alpha\in[0.11,\,0.29]$에 있다.
$\alpha<10^{-2}$인 $54$번은 $50$ epoch 동안 학습이 느려서 최고 $77.6\%$에 그쳤고,
$\alpha\ge0.3$인 $12$번은 평균 $29\%$로 대부분 발산하였다.
$[0.05,\,0.3)$ 사이에서도 $\alpha>0.25$인 $4$번 중 $2$번은 실패하였다 ($10.1\%$, $62.0\%$).
\item 가중치 감소: 상위 $10$개의 $\lambda$가 $10^{-6}$에서 $4\times10^{-2}$까지 고르게 퍼져 있다.
$\alpha\in[0.05,0.3)$인 $20$번만 보아도 $\lambda$와 검증 정확도 사이에 뚜렷한 관계가 없다.
\end{itemize}

\textbf{좁힌 범위.} 학습률은 $\boxed{\alpha\in[0.05,\,0.25]}$ (로그 척도)로 좁혔다.
가중치 감소는 이 실험에서 성능을 거의 좌우하지 않으므로 범위를 좁힐 근거가 없어서 $\lambda\in[10^{-6},\,10^{-1}]$ (로그 척도)을 그대로 두었다.
(근거 없이 좁히면 좋은 값을 놓칠 수 있다. 특정 $\lambda$ 근처를 더 보고 싶다면 학습률을 고정하고 $\lambda$만 탐색하는 편이 낫다.)

\textbf{결과 ($30$번).} 나머지 조건은 원래와 같다 (훈련 이미지 $500$장, 은닉층 $[100]\times6$, He, SGD, $50$ epoch).
\begin{center}
\begin{tabular}{lcc}
\toprule
 & 원래 ($100$번) & 좁힌 범위 ($30$번) \\
\midrule
학습률 $\alpha$의 범위 & $[10^{-5},\,1]$ & $[0.05,\,0.25]$ \\
가중치 감소 $\lambda$의 범위 & $[10^{-6},\,10^{-1}]$ & $[10^{-6},\,10^{-1}]$ \\
가장 좋은 검증 정확도 & 0.8714 & 0.8710 \\
검증 정확도의 평균 & 0.4602 & 0.8398 \\
검증 정확도 $50\%$ 미만 & $54$번 & $1$번 \\
고른 $\alpha$, $\lambda$ & $0.174$, $5.3\times10^{-6}$ & $0.0939$, $2.8\times10^{-2}$ \\
그 network의 시험 정확도 & 0.8701 & 0.8714 \\
\bottomrule
\end{tabular}
\end{center}
좁힌 범위의 상위 $5$개: $(\alpha,\lambda,\text{검증})=(0.0939,\ 0.0283,\ 0.8710)$, $(0.109,\ 0.0280,\ 0.8668)$,
$(0.103,\ 0.0096,\ 0.8660)$, $(0.151,\ 0.0018,\ 0.8646)$, $(0.125,\ 4.1\times10^{-5},\ 0.8646)$.

\textbf{해석.} 범위를 좁히자 $30$번 중 $29$번이 $80\%$ 이상이었고, 시도의 평균이 $46\%$에서 $84\%$로 올랐다.
즉 같은 계산량을 쓸모없는 영역에 낭비하지 않게 되었다.
하지만 가장 좋은 검증 정확도는 $87.1\%$로 원래와 거의 같다.
이미 원래 탐색이 좋은 영역을 충분히 찾았고, 이 network와 훈련 이미지 $500$장으로 얻을 수 있는 정확도가 이 정도라는 뜻이다.
정확도를 더 올리려면 하이퍼파라미터보다 훈련 데이터의 양이나 network의 구조를 바꾸어야 한다.
(시험 정확도는 고른 뒤 한 번만 계산하였다.)
\end{sol}

\end{problems}

\end{document}
